%%%%%%%%%%%%%%%%%%%%%%%%%%%%%%%%%%%%%%%%%%%%%%%%%%%%%%%%%%%%%%%%%%%%%%
% 中国A股市场波动率管理投资组合(VMP)的实证研究
% ——基于Fama-French五因子与动量因子的检验
%
% 基于 njuthesis 模板
%%%%%%%%%%%%%%%%%%%%%%%%%%%%%%%%%%%%%%%%%%%%%%%%%%%%%%%%%%%%%%%%%%%%%%

\documentclass[
    % type = master,
    % degree = academic,
    decl-page,
    oneside,
  ]{njuthesis}

\input{njuthesis-setup.def}

% 载入所需宏包
\usepackage{booktabs}    % 专业三线表
\usepackage{multirow}    % 表格合并行
\usepackage{amsmath}     % 数学公式
\usepackage{graphicx}    % 插图
\usepackage{float}       % 强制浮动体位置
\usepackage{listings}    % 代码排版
\usepackage{xcolor}      % 颜色支持

% 缩小章标题与正文之间的间距
\ctexset{
  chapter/afterskip = 20pt
}

% Python代码样式设置
\lstdefinestyle{pythonstyle}{
  language=Python,
  basicstyle=\small\ttfamily,
  keywordstyle=\color{blue}\bfseries,
  commentstyle=\color{gray}\itshape,
  stringstyle=\color{red!70!black},
  numbers=left,
  numberstyle=\tiny\color{gray},
  numbersep=8pt,
  breaklines=true,
  frame=single,
  rulecolor=\color{black!30},
  backgroundcolor=\color{gray!5},
  showstringspaces=false,
  tabsize=4,
  captionpos=b,
  morekeywords={as, True, False, None, with, lambda},
}

% 开始编写论文
\begin{document}

%---------------------------------------------------------------------
%	封面、摘要、前言和目录
%---------------------------------------------------------------------

\maketitle
\raggedbottom

% ====================================================================
% 中文摘要
% ====================================================================
\begin{abstract}

因子投资在实践中常伴随收益的剧烈波动乃至“因子崩盘”等极端风险。2017年，Moreira和Muir提出波动率管理投资组合（VMP）策略，以已实现方差的倒数动态调整因子仓位，这一策略在美股市场取得了显著的风险调整收益提升，但在中国A股市场的有效性尚缺乏系统检验。

本文基于CSMAR数据库2005—2025年的日度数据，将VMP策略分别应用于中国A股的Fama-French五因子（MKT、SMB、HML、RMW、CMA）及动量因子（UMD）。为回应DeMiguel等人（2020）关于交易成本的质疑，本文施加了3倍最大杠杆约束，并按每日权重变化幅度扣除单边万分之十的交易成本，统计推断采用Newey-West HAC标准误。

实证分析结果显示，VMP策略对不同因子的影响并不一致。市场因子和规模因子经VMP管理后绩效不升反降，其中规模因子的Alpha显著为负（$-5.4\%$，$p=0.018$）。价值、盈利和投资三因子虽有改善但并不显著。而动量因子经过VMP管理后的增效显著，具体表现为：UMD年化收益率由$-1.9\%$翻转为$+6.8\%$，净夏普比率由$-0.11$上升至$0.44$，最大回撤由$77.1\%$降至$55.9\%$，Spanning回归的年化Alpha达到$8.1\%$。进一步的稳健性检验表明，在63天和126天窗口下，UMD 的 alpha 仍在$1\%$水平显著；在2015年股灾和2018年贸易战期间，VMP 也表现出一定的回撤缓释作用。

总体来看，VMP的改善作用主要集中在动量因子上。这说明，A股动量效应的弱表现未必意味着动量信号本身不存在，更可能是高波动时期的极端损失掩盖了原有的收益信号。

\end{abstract}

% ====================================================================
% 英文摘要
% ====================================================================
\begin{abstract*}

Can volatility-managed portfolios (VMP)---which scale factor exposures inversely with recent realized variance---improve risk-adjusted returns in China's A-share market? The answer, it turns out, depends heavily on which factor one applies the strategy to.

Using CSMAR daily data Spanning 2005--2025, this paper constructs VMP versions of the five Fama-French factors plus momentum (UMD) under realistic constraints: a 3x leverage cap and one-sided transaction costs of 10 basis points on daily weight turnover. Newey-West HAC standard errors are employed throughout.

The key finding is a striking asymmetry. For momentum, VMP works remarkably well---flipping UMD's annualized return from $-1.9\%$ to $+6.8\%$ and pushing the net Sharpe ratio to $0.44$. The Spanning-regression alpha of $8.1\%$ per year ($t = 3.27$) remains significant at the $1\%$ level across alternative volatility windows (63-day, 126-day) and during both the 2015 crash and 2018 trade-war episodes. For market and size factors, however, VMP backfires: size-factor alpha is $-5.4\%$ ($p = 0.018$), since these factors earn their premiums precisely during high-volatility rebounds that VMP scales away from. Value, profitability, and investment factors show near-zero alphas either way.

These results echo---from an independent emerging-market sample---the conjecture of Cederburg et al.\ (2020) that VMP's benefits concentrate on momentum-type strategies. They also suggest that the apparent ``failure'' of momentum in A-shares may stem not from an absence of signal, but from a few catastrophic tail events that mask it.

\end{abstract*}

% 生成目录
\tableofcontents
\listoffigures
\listoftables

%---------------------------------------------------------------------
%	正文部分
%---------------------------------------------------------------------
\mainmatter

% ====================================================================
% 第一章：引言
% ====================================================================
\chapter{引言}

\section{研究背景与研究意义}

\subsection{研究背景}

因子投资是现代资产定价理论的核心范式。自Fama和French \cite{fama1993common} 提出三因子模型以来，学术界陆续发现了包括市场（MKT）、规模（SMB）、价值（HML）、盈利（RMW）、投资（CMA）和动量（UMD）在内的多种定价因子 \cite{fama2015five,carhart1997persistence,jegadeesh1993returns}，这些因子被广泛应用于投资组合构建和业绩归因。然而，因子投资的收益并不稳定——因子组合往往伴随着剧烈的波动率波动，尤其在市场极端行情中，因子收益可能出现大幅回撤。动量因子的“动量崩盘”（Momentum Crash）现象便是一个典型案例：Daniel和Moskowitz \cite{daniel2016momentum} 发现，动量策略在少数极端月份的巨额亏损可以吞噬数年的累积收益。

如何在不放弃因子溢价的情况下，对因子组合的波动率风险进行较有效管理，已经成为近年来学术界以及业界比较关注的问题。Moreira和Muir \cite{moreira2017volatility} 提出了波动率管理组合（Volatility Managed Portfolio，VMP）策略.其核心思想即利用已实现波动率具有一定时序可预测性这一特点，在市场波动率较高时期减少因子暴露，而在波动率较低时期增加因子暴露。其研究结果表明，在美国市场主要因子中应用 VMP策略之后，所得到的管理组合相比原始因子具有显著为正的 Alpha，也就是说，VMP组合能够在一定程度上提供超过原始因子的风险调整后收益。该研究提出之后，也在学术界引起较多讨论。

不过，VMP策略是否具有普遍有效性，学界至今并未形成一致结论。Cederburg等人 \cite{cederburg2020performance}从实证角度指出，Moreira和Muir所报告的alpha虽然在Spanning test框架下具有统计显著性，但是如果进一步考虑实际交易限制条件，其对应经济收益并不是非常明显，而且该策略的有效性基本只集中在动量类因子；第二是在理论层面，Guo和Li (2019) 认为Moreira和Muir在推导VMP组合Alpha时，使用的收益率度量方法并不合适，在进行修正之后，无法从理论上证明VMP组合一定存在显著为正的Alpha。此外，Liu等人 \cite{liu2019size} 和Chi等人 (2019) 的研究进一步表明，VMP策略在不同市场环境以及制度背景之下，其适用性也存在比较明显差异。

在这一背景下，中国A股市场为重新检验VMP策略提供了一个合适的研究对象。与美国市场相比，A股在投资者结构、交易制度和因子表现上都具有较强的特殊性。A股市场以散户投资者为主，交易行为更容易受到情绪驱动，因而波动率的形成机制与成熟市场并不完全相同。与此同时，中国市场的因子定价结构也与美国存在差异，Liu等人 \cite{liu2019size} 指出，Fama-French因子在A股的适用性有限，尤其动量效应往往表现较弱甚至呈现反转。此外，A股市场还有涨跌停制度、T+1交易约束以及政策冲击频繁等特点。Chi等人 (2019) 指出，A股收益率与历史波动率之间呈正相关关系，这与美国市场的负相关特征形成鲜明对比，也意味着 VMP 策略在中国市场未必会重复其在美国市场中的表现。

在此背景下，一个学术问题自然浮现：\textbf{波动率管理策略在中国A股市场是否有效？如果有效，对哪些因子有效？其有效性的来源和机制是什么？}现有文献对此缺乏系统性的实证检验。本文正是在此背景下，对中国A股市场六大主要因子（MKT、SMB、HML、RMW、CMA、UMD）的VMP策略进行全面的实证评估。

\subsection{研究意义}

Moreira和Muir \cite{moreira2017volatility} 的原始研究基于美国市场数据，其结论能否推广至新兴市场尚缺乏充分验证。中国A股是全球第二大股票市场、最大的新兴市场，其制度环境、投资者结构同美国市场存在着根本性的差别。本文对VMP策略在A股六个因子上的系统检验，可以回答VMP的有效性是否具有跨市场普适性这个问题。与此同时，Cederburg等人 \cite{cederburg2020performance} 基于美股数据提出VMP只对动量类因子显著的结论与本文基于中国股市市场的实证结果中具有一致性，这就给这一结论提供了一个独立的跨市场验证。更进一步来说，目前的文献大多认为中国A股不存在动量效应\cite{liu2019size}，本文发现波动率管理策略可以使得动量因子产生年化$8.1\%$的显著Alpha，这就说明A股动量效应的“失效”并不是由动量信号本身的缺失所造成的，而是由于少数高波动的极端事件对动量收益造成了毁灭性的冲击，VMP在高波动时期自动降低了仓位过滤掉了这些极端噪音，使潜在的动量信号得以显现出来。这一发现给认识中国市场的动量异象赋予了波动率方面的新视角。

除此之外，本文对于量化投资实践也有一定的借鉴意义。从研究结果可知，在A股市场中已经具有较好风险收益特征的因子（市场因子、规模因子）如果盲目地使用VMP策略，不但不能提高绩效，反而会严重损害投资效果。这就提醒量化投资管理者，在使用波动率管理策略的时候要慎重地区分出因子的特征。就动量策略而言，虽然裸动量在中国市场中表现不佳，但是结合波动率管理之后，其夏普比率由$−0.11$提升至$0.44$。这意味着，经由恰当的风险管理框架，原本被认为在中国市场无效的动量策略，可以被转化为具备投资价值的工具。本文对于2015年股灾、2018年贸易战等极端时期VMP策略差异化的分析，给投资者在不同的危机环境下进行风险管理决策提供经验。

\section{研究思路与研究方法}

本文的研究主要分为三个步骤。首先，对VMP策略的相关文献进行梳理，说明其理论基础、核心机制以及现有争议。其次，基于中国A股市场六个因子的日度数据分别构建VMP组合，并通过Spanning回归检验其是否能够带来超出原始因子的增量收益。最后，通过替换波动率窗口、划分样本区间以及考察极端市场时期的表现，对前述结论进行稳健性检验。

研究方法方面，波动率估计采用滚动窗口已实现方差（基准为21个交易日），作为因子条件方差的实时度量。策略权重按照Moreira和Muir \cite{moreira2017volatility} 的方法以方差倒数确定，同时施加3倍最大杠杆约束来确保策略在实践中可行，并按每日权重变化幅度扣除单边万分之十的交易成本——这一处理主要是回应DeMiguel等人 \cite{demiguel2024transaction} 关于交易成本侵蚀VMP收益的批评。检验环节则使用Spanning回归框架：把VMP管理后的因子净收益对原始因子收益做回归，看截距项（Alpha）是否显著异于零即可。统计推断一律使用Newey-West异方差自相关一致标准误，以避免日度数据中常见的条件异方差和短期序列相关导致的推断失真。

\section{研究内容与章节安排}

本文以Fama-French五因子（MKT、SMB、HML、RMW、CMA）加动量因子（UMD）共6个核心因子为研究对象，使用CSMAR数据库2005年至2025年的日度数据，在日度频率下构建VMP组合并进行系统的实证评估。全文共六章。本章为引言；第二章梳理因子投资理论、VMP策略争论及A股因子有效性方面的已有文献；第三章从各类因子的“波动率—收益”时序结构出发推导三个核心假设；第四章介绍数据来源、策略构建方法以及Spanning回归检验框架；第五章报告描述性统计、绩效对比、Alpha检验和多维度稳健性分析的结果；第六章总结结论并讨论理论贡献、实践启示和研究局限。

\section{可能的创新}

本文的可能创新主要体现在两个方面。第一，在研究对象上，现有关于VMP的研究多集中于单一因子，或者主要采用月度频率展开分析。本文则在日度频率下同时考察Fama-French五因子与动量因子，并在杠杆约束和交易成本扣除的条件下评价策略表现，更完整地考察VMP在中国A股核心因子上的适用性。第二，在研究发现上，本文尝试从波动率管理的角度重新理解A股市场中动量因子的弱表现。已有研究通常认为中国市场的传统动量效应并不显著，而本文发现，经VMP处理后，动量因子能够获得显著为正的alpha。这说明，A股动量信号的弱表现未必完全意味着信号本身不存在，也可能是高波动时期的极端反转损失掩盖了其原有收益。就这一点而言，本文的结果也为Cederburg等人 \cite{cederburg2020performance} 关于“VMP仅对动量类因子稳健”的判断提供了来自中国市场的补充证据。

% ====================================================================
% 第二章：文献综述
% ====================================================================
\chapter{文献综述}

\section{因子投资理论}

\subsection{从CAPM到Fama-French五因子模型}

因子投资理论最早可以追溯到Sharpe \cite{sharpe1964capital} 提出的资本资产定价模型（CAPM）。CAPM认为，市场投资组合是唯一的风险定价因子，个股的预期超额收益主要由其市场Beta决定。但后续大量实证研究发现，CAPM对股票收益截面差异的解释能力是有限的。特别是小市值股票和高账面市值比股票的平均收益，往往明显高于CAPM模型的预测值，这说明仅依靠市场因子并不足以解释股票收益的系统性差异。

在这一背景下，Fama和French \cite{fama1993common} 提出了三因子模型，在市场因子之外加入了规模因子（SMB）和价值因子（HML）。其中，SMB反映小市值股票相对于大市值股票的超额收益，HML反映高账面市值比股票相对于低账面市值比股票的超额收益。三因子模型显著提高了对股票截面收益率的解释能力，也因此成为资产定价研究中的经典框架。围绕规模效应和价值效应的来源，学界一直存在两种主要解释：一种认为这类收益来源于对更高风险的补偿，另一种则认为其背后包含投资者错误定价和行为偏差的影响。

不过，三因子模型仍然无法覆盖所有重要的定价异象。Novy-Marx \cite{novymarx2013other}发现，在控制账面市值比之后，盈利能力更强的公司仍然具有更高的平均收益。与此同时，也有研究指出，投资行为较为稳健的公司往往比激进扩张的公司表现更好。基于这些经验事实，Fama和French进一步提出五因子模型，在原有三因子的基础上加入盈利因子（RMW）和投资因子（CMA）。五因子模型的提出拓展了经典因子框架，但也引发了新的讨论。例如，Fama和French自己就指出，在纳入RMW和CMA之后，HML的边际解释力明显下降，这使得价值因子是否冗余成为后续研究中的一个重要问题。

\subsection{动量效应与动量崩盘}

Jegadeesh和Titman \cite{jegadeesh1993returns} 最早系统性地记录了股票市场中的动量效应：买入过去3至12个月表现最好的股票（赢家组合），同时卖空表现最差的股票（输家组合），可以在未来获得显著的正收益。这一发现对有效市场假说构成了重大挑战，因为历史价格信息在该假说下不应能够预测未来收益。Carhart \cite{carhart1997persistence} 在Fama-French三因子的基础上加入动量因子（UMD），形成了四因子模型，动量因子由此成为资产定价研究中的标准组成部分。

然而，动量策略或存在较为严重的尾部风险。Daniel和Moskowitz \cite{daniel2016momentum} 对动量崩盘（Momentum Crashes）现象展开了系统性的研究。其研究发现，在市场由熊市向牛市急剧反转的初期，动量策略遭受具有毁灭性质的亏损。其作用机制在于：在经历熊市之后，输家组合中的股票往往具备较高的市场Beta，当市场出现突然的反弹时，此类高Beta之输家股票的涨幅或将远超赢家股票，进而导致多空两端同时遭受损失。Daniel与Moskowitz的数据表明，在1932年美股大萧条反弹期间，动量策略于两个月内的亏损幅度超过了$90\%$。此种极端的尾部风险，使得动量策略的无条件夏普比率被大幅拉低，这同时亦促使学术界开始探索如何通过风险管理手段以改善动量策略的收益分布特征。

此外，顾明等 \cite{guming2024factor} 对中国A股市场的因子动量效应进行了系统检验。所谓因子动量，是指在多个已知因子之间运用动量策略进行配置，即增配近期表现好的因子、减配表现差的因子。他们发现时序因子动量和截面因子动量策略在A股市场均能盈利，月均超额收益率分别为0.53\%和0.54\%，且在2008年、2015年等极端行情中未出现崩溃。需要指出的是，因子动量与本文研究的个股动量因子（UMD）属于不同层面的概念：前者是因子之间的配置策略，后者是基于个股历史收益排序构建的单一因子。顾明等的发现表明，中国市场中因子收益本身存在时序延续性，这为理解波动率管理策略的作用机制提供了有益的背景信息。

\section{波动率管理投资组合（VMP）理论}

\subsection{VMP的提出与核心思想}

关于波动率择时是否具有经济价值，Fleming等人 \cite{fleming2001economic} 给出了最早的系统性证据。他们观察到已实现波动率在时间序列上有较强的可预测性，据此动态调整组合权重可以给投资者带来年化50到200个基点的确定性等价收益提升。这个结果为后来Moreira和Muir的工作奠定了经验上的基础。

Moreira和Muir \cite{moreira2017volatility} 在此基础上把波动率择时的思想发展为一个完整的框架。他们抓住了一个关键的经验事实：因子收益的条件方差在时间序列上高度可预测（所谓波动率聚集），但条件均值跟条件方差之间的关系很弱，甚至可能是负的。这意味着波动率升高时投资者多承担了风险却没得到相应补偿——那么反过来操作，高波动时减仓、低波动时加仓（用方差倒数缩放仓位），就能在不太牺牲期望收益的前提下把波动率降下来，夏普比率自然就上去了。他们用美股多个因子做了实证，发现VMP在大部分因子上都能产生显著的Spanning回归alpha。

Barroso和Santa-Clara \cite{barroso2015momentum} 则专门从动量策略的角度给出了更直接的支持。动量策略的已实现波动率有很强的一阶自相关（系数超过0.9），利用这一点做实时风险缩放，能将夏普比率从0.53拉到0.97。更关键的是，经过风险管理的动量策略几乎完全消除了崩盘式亏损：2009年3月美股大反弹时，裸动量亏了超过70\%，管理后的版本只亏了不到10\%。这意味着波动率管理对高尾部风险因子有特别突出的改善效果。

\subsection{对VMP的质疑与学术争论}

有关VMP策略是否有效，学术界一直存在争议。现有批评主要集中在两个方面。其一是统计稳健性问题。Cederburg等人 \cite{cederburg2020performance}对Moreira和Muir的研究结论进行了重新检验，发现VMP的改善效果在很大程度上依赖于因子类型和检验设定。具体来看，VMP对市场因子（MKT）的改善相对稳健，但对规模因子（SMB）、价值因子（HML）等其他因子的作用较弱，在大多数情况下并不显著。此外，VMP的表现对波动率估计窗口也较为敏感，使用21天、63天、126天等不同回望窗口时，所得结果并不一致。因此，他们认为Moreira和Muir原始结论中的部分显著性，可能与参数选择偏差有关。其二是交易成本问题。DeMiguel等人 \cite{demiguel2024transaction}从实践可行性的角度指出，VMP需要根据波动率变化频繁调整仓位，在日度调仓设定下会产生较高换手率。以市场因子为例，VMP的年化换手率可能明显高于原始买入持有策略；对于流动性较差的因子，如以小市值股票为基础构建的SMB，交易成本的影响则更为明显。他们的估算结果表明，在纳入较为合理的交易成本假设后，多数因子的VMP收益改善将不再显著，部分因子甚至会出现净亏损。

上述争论说明，VMP的有效性并不能被简单视为一个已经得到定论的问题。支持者强调波动率管理对风险收益特征的改善作用，批评者则认为这种改善在考虑参数敏感性和实际交易摩擦后可能难以稳定成立。近年的研究并没有消除这一分歧，反而进一步表明，VMP的表现会随着研究框架和市场环境的变化而发生差异。例如，DeMiguel等人\cite{demiguel2024multifactor}在多因子框架下发现，波动率管理带来的收益改善明显弱于单因子检验中的结果，说明交易成本和因子拥挤可能限制其实际表现；Schwarz\cite{schwarz2025international}也发现，VMP在不同国家市场中的效果并不一致，在新兴市场中的表现尤其不稳定。基于这一背景，本文在实证分析中不仅考察VMP策略是否能够改善因子收益表现，同时也扣除了交易成本，以更全面地评估其在A股市场中的适用性。

\section{中国A股市场因子有效性}

\subsection{中国市场的因子格局}

中国A股市场的因子结构同成熟市场相比，存在着比较明显的差别，近些年来已有不少研究对这一现象做了系统的检验。Liu等人 \cite{liu2019size} 针对中国市场的特殊性，构建了中国三因子模型（CH-3），其关键的创新之处在于：将流通市值最小的$30\%$股票识别为受壳价值污染的样本并加以剔除，剔除的原因是这些微型股的高收益率主要反映的是借壳上市产生的投机价值，而非真实的规模溢价。研究还发现，在剔除壳污染之后，中国市场仍然存在显著的规模效应，但是账面市值比（B/M）的解释力却大幅度降低，而市盈率倒数（E/P）却反映出了更强的定价能力。
在中文学术界，赵胜民等 \cite{zhao2016volatility} 较早对Fama-French五因子模型在A股市场的适用性进行了系统检验，发现五因子模型相对于三因子模型的改进幅度有限，盈利因子和投资因子在中国市场的解释力不如在美股中表现得那样突出。李志冰等 \cite{li2017factor} 进一步利用2000—2015年A股数据对五因子模型进行了全面实证检验，发现规模因子和价值因子在中国市场仍然具有显著的定价能力，但五因子模型对于特定类型异象（如应计异象）的解释力有所不足。

Hu等人 \cite{hu2019chinese} 也对中国市场的Fama-French因子进行了较为全面的检验。其结果表明，在各类因子中，规模因子（SMB）对中国股票收益截面差异的解释力最强，无论是在时间序列回归中，还是在Fama-MacBeth截面回归中，都具有较高的显著性。相比之下，价值因子（HML）的表现不够稳定，在不同的检验设定和样本期间下，结论存在较大差异。Blitz等人 \cite{blitz2021volatility}则从低波动率异象的角度揭示了中国市场的另一项特征。他们发现，A股市场中低风险股票的收益率会系统性地高于高风险股票，而且这一现象主要由波动率驱动，而不是由Beta驱动。Liu等人 \cite{liu2024trend}从趋势因子的角度进一步说明了中国市场因子结构的特殊性。他们发现，传统动量在A股中的表现并不理想。

总体来看，已有研究说明，中国市场中的规模、价值、波动率和趋势等信息都表现出不同于成熟市场的定价特征。这些研究一方面表明，国际经典因子模型在A股市场中的适用性需要结合本土市场特征重新理解；另一方面也为后文分析VMP策略在A股市场中的作用提供了经验背景。

\subsection{中国的动量困境}

在中国市场的各类因子中，动量因子的表现尤为特殊。与美股市场中较为稳健的动量效应不同，大量实证研究发现，传统个股动量在A股市场并不显著，甚至在部分样本中表现为反转。Cheema和Nartea \cite{cheema2018momentum}对1996年至2013年的A股数据进行检验后发现，基于过去6个月收益率排序构建的动量策略会产生显著为负的收益，这与美股市场中的强动量效应形成了明显反差。

对于这一现象，现有文献给出了几种可能的解释。首先，中国市场长期以散户投资者为主，散户交易频繁且更容易出现过度反应，这使得短期价格更容易出现反转，而不是持续延续。其次，涨跌停板制度在一定程度上限制了极端信息的即时反映，影响了价格发现过程。再次，IPO暂停、杠杆管控等政策冲击容易带来市场风格的突然切换，从而削弱动量效应的持续性。也就是说，中国市场中的动量失效并不一定意味着价格中完全不存在趋势信息，更可能说明这种信息在个股层面容易受到制度环境和交易行为的干扰。

不过，个股动量的弱表现并不意味着动量类信息完全没有价值。已有研究发现，在更高的聚合层面上，动量效应仍然可能存在。顾明等 \cite{guming2024factor}发现，A股市场中的因子动量效应是显著的，其基于10个非动量因子构建的因子动量策略年化收益率达到$9.91\%$，并且在经过多因子模型调整之后仍然显著。Ma等人 \cite{jiang2024factormomentum} 进一步检验了中国A股中的时序因子动量和截面因子动量，发现两类策略在不同设定下均能够获得显著的超额收益。这些结果表明，在因子层面动量信息仍具有定价价值。一个可能的解释是：个股层面的高特质性噪音削弱了动量信号，而当组合聚合到因子层面时，这种信号更容易显现。

\subsection{VMP在中国市场的研究现状}

VMP策略在中国A股市场上的适用性近些年来越来越受到学术界的重视，目前已有两项与其直接相关的研究成果。Wang和Li \cite{wang2024vmp} 利用2000年1月至2022年6月窗口内A股数据，构造出71个特征管理因子的VMP版本，在月度频率下进行了系统的检验。其发现VMP在样本内、样本外都比原始因子更优，优势主要来自于波动率缩放策略在高波动时期降低仓位暴露造成的。

Zhong和Li \cite{zhong2024sentiment} 则从投资者情绪的角度，探讨了中国市场中VMP超额收益的来源。其构建了理论模型以刻画投资者情绪对VMP风险调整后收益的影响机制，并利用A股数据加以实证检验。研究发现，SMB因子经过波动率管理之后，夏普比率出现了显著的提升，且此种改善在低情绪时期尤为明显。此一结果表明，中国市场由散户主导、受情绪驱动的投资者行为特征或许恰恰是VMP策略得以发挥作用的条件之一。Guo等人 \cite{guo2024tailrisk} 从尾部风险的视角进一步拓展了VMP研究的分析框架，发现波动率管理策略的有效性与因子收益分布的尾部特征密切相关：对于左尾风险较高的因子，VMP能够通过在极端市场状态下降低暴露来系统性地改善其风险调整收益，这一发现为理解VMP在中国市场的差异化表现提供了新的理论视角。

\section{本章小结}

以上从三个方向梳理了与本文相关的已有文献。因子投资理论从CAPM一路发展到五因子模型，对截面收益差异的解释力不断提升，但动量效应在中国市场的失效仍是一个悬而未决的问题。VMP方面，Moreira和Muir开创性地证明了波动率择时对因子收益的改善潜力，但Cederburg等人和DeMiguel等人分别从统计稳健性和交易成本两个角度提出了有力质疑。就中国市场而言，因子结构与成熟市场差别明显，VMP在A股的研究仍比较初步。

综合来看，一个明显的空白在于：对A股核心因子施加VMP后，在日度频率、有交易成本、有极端行情的条件下，到底能不能带来可靠的风险调整收益改善？本文试图填补该研究空白。


% ====================================================================
% 第三章：研究假设
% ====================================================================
\chapter{研究假设}

前文从因子投资理论、VMP策略及其争论、以及中国A股市场因子有效性三个方面对相关文献做了梳理。本章在此基础上，针对VMP策略作用于不同类型因子时的效果，分别进行理论推导并提出研究假设。

\section{VMP策略对动量因子的改善效果}

波动率管理组合（VMP）策略的基本操作是以因子已实现方差的倒数作为当期仓位权重。当某因子近期的波动率处于高位时，VMP自动缩减该因子的暴露；当波动率回落至低位时，则相应增大暴露。Moreira和Muir \cite{moreira2017volatility} 指出，这一操作能够带来风险调整收益改善的前提条件是：因子的条件期望收益在高波动时期没有同比例增大。换言之，只要高波动时期因子的“性价比”（单位风险的收益补偿）有所下降，VMP通过减少这些低性价比时段的暴露，就能提升组合整体的夏普比率。

动量因子恰好满足这一前提条件，而且满足得非常彻底。Daniel和Moskowitz \cite{daniel2016momentum} 系统记录了“动量崩盘”现象：动量策略在绝大多数交易日里贡献稳定的小幅正收益，但在少数市场急剧反转的时段中，会出现毁灭性的集中亏损，数日的巨额亏损足以抵消数年的累积收益。Barroso和Santa-Clara \cite{barroso2015momentum} 进一步发现了一个关键的规律性事实——动量崩盘几乎总是发生在事前波动率处于极端高位的时候，动量策略的已实现波动率对未来崩盘具有很强的预测力。也就是说，对于动量因子而言，高波动不意味着高收益，反而意味着高崩盘风险。其条件期望收益在高波动时期是负的，在正常波动时期才是正的。

中国A股市场的若干特征使得上述机制更加突出。如文献综述中已经提到的，A股个股动量在无条件均值检验中表现不显著甚至为负。但如果拆分动量策略的收益分布来看，把整体表现拖到负值区域的，主要是2015年杠杆牛市崩盘、2018年贸易战引发的持续下跌、以及2022年疫情封控这几段高波动时期里出现的集中巨亏。这些极端反转事件发生的前期，动量组合的波动率都有比较明显的上升。A股市场散户主导、政策冲击频繁的特点使得这类极端事件比美股更加集中和剧烈。Cederburg等人 \cite{cederburg2020performance} 在美股数据上已经发现VMP主要对动量因子有效这一规律；在A股市场中，由于动量崩盘事件更为极端，VMP通过事前波动率信号降低暴露所能过滤掉的尾部损失也更大，被过滤掉这些极端亏损之后，动量信号中潜在的正收益就会得以显现。

综合以上分析，本文提出：

\textbf{假设一：VMP策略能够显著改善动量因子在中国A股市场的风险调整收益。}

\section{VMP策略对市场因子和规模因子的损害效果}

VMP策略能否改善某个因子，取决于该因子自身的“波动率—收益”时序结构。如上一节所述，VMP发挥正面作用的前提是高波动时期因子的条件期望收益没有同比例升高。如果这个前提不成立——也就是某个因子在高波动时期反而能够提供更高的期望收益，波动率和预期溢价之间呈正相关——那么VMP在高波动期的减仓操作就会系统性地把那些高回报的交易日给剔除出去，结果组合的表现反而不如原来的等权暴露。

中国A股的市场超额收益因子和规模因子恰好呈现出这样的结构。Chi等人 (2019) 发现了一个与美股截然不同的现象：A股市场中，历史波动率与未来收益率之间呈现正相关关系，高波动时期往往对应着更高的预期超额收益。这一点在行为金融学框架下可以这样理解——A股以散户投资者为主体，高波动时期伴随恐慌性抛售，容易形成过度折价，后续的均值回复则表现为正向收益补偿。规模因子（SMB）的情况也比较类似：小盘股本身波动率就高于大盘股，而小盘溢价在市场剧烈波动的阶段——比如流动性恐慌、政策急转弯等——反而容易出现集中释放。

在这样的“高波动—高溢价”结构之下，VMP在波动率攀升时降低仓位暴露，实质上是把因子溢价集中兑现的那些交易日给低权重了。减仓的时候正是因子赚钱最多的时候，加仓的时候对应的反而是低溢价的平静期。这种操作方向与因子收益方向之间的背离不是偶然的参数敏感问题，而是由A股市场因子和规模因子“高波动—高溢价”这一结构性特征所内在决定的。只要这一结构性特征存在，VMP的降仓操作就必然系统性地损害这两个因子的收益获取能力。

综合以上分析，本文提出：

\textbf{假设二：VMP策略显著损害市场因子和规模因子在中国A股市场的绩效表现。}

\section{VMP策略对价值因子、盈利因子和投资因子的中性效果}

前两节分别分析了VMP策略在两种极端的“波动率—收益”结构下的作用方向：动量因子的高波动对应高崩盘风险，VMP通过减仓过滤尾部损失而产生正面效果；市场因子和规模因子的高波动对应高溢价释放，VMP通过减仓错过收益高峰而产生负面效果。这两种机制能否发挥作用，关键在于因子的条件期望收益与其自身波动率之间是否存在足够强的、方向明确的时序相关关系。

价值因子（HML）、盈利因子（RMW）和投资因子（CMA）在这方面的特征与前述两类因子有所不同。从已有文献来看，这三个因子的溢价主要来源于横截面定价——价值股相对成长股的估值修复、高盈利企业相对低盈利企业的盈利持续性、以及保守投资企业相对激进投资企业的资本效率差异。这些横截面收益差异的驱动力量是企业基本面特征，而非市场整体的波动率状态。与市场因子和规模因子不同，现有研究并未发现这三个因子在高波动时期系统性地释放更高溢价；与动量因子不同，这三个因子也不存在与波动率状态高度相关的崩盘风险模式。

换言之，价值因子、盈利因子和投资因子的条件期望收益与其自身波动率之间缺乏足够强的时序联动——既不是正相关（高波动=高收益），也不是负相关（高波动=高亏损），而是近似无关。在这种“波动率中性”的结构下，VMP在高波动期的减仓操作既不会像作用于动量因子时那样过滤掉集中亏损（因为没有系统性的高波动崩盘），也不会像作用于市场因子时那样错过集中溢价（因为没有系统性的高波动高收益），其净效果预期接近于零。

综合以上分析，本文提出：

\textbf{假设三：VMP策略对价值因子、盈利因子和投资因子不产生显著影响。}

% ====================================================================
% 第四章：研究设计
% ====================================================================
\chapter{研究设计}

\section{数据描述}

\subsection{数据来源}

本文所用到的因子收益率数据来自于国泰安（CSMAR）数据库中的“Fama-French因子”。本文用中国A股市场Fama-French五因子日度收益率来表示，即市场超额收益因子（$\text{MKT\_RF}$）、规模因子（$\text{SMB}$）、价值因子（$\text{HML}$）、盈利因子（$\text{RMW}$）、投资因子（$\text{CMA}$）。本文还从CSMAR数据库中得到了动量因子（$\text{UMD}$）的日度收益率。无风险利率（$\text{RF}$）采用CSMAR提供的日度化一年期定期存款利率。

以下对六个因子的构建方法进行简要说明：市场超额收益因子（MKT\_RF）定义为A股市场全部股票的流通市值加权收益率减去无风险利率。规模因子（SMB）和价值因子（HML）的构建遵循Fama和French \cite{fama1993common} 的经典方法：每年6月末，按流通市值中位数将全部股票分为小盘（S）和大盘（B）两组，同时按账面市值比（B/M）的30\%和70\%分位数将股票分为高（H）、中（M）、低（L）三组，由此形成$2 \times 3$共六个投资组合。SMB为三个小盘组合与三个大盘组合的等权收益率之差，HML为两个高B/M组合与两个低B/M组合的等权收益率之差。盈利因子（RMW）以营业利润率为排序变量，投资因子（CMA）以总资产增长率为排序变量，构建方式与HML类似，分别取高盈利/低盈利之差和保守投资/激进投资之差。动量因子（UMD）依据过去2至12个月（跳过最近1个月）的个股累计收益率进行排序，做多过去赢家组合、做空过去输家组合。所有因子均采用流通市值加权方式计算组合收益率。

CSMAR数据库是中国金融学实证研究中使用最为广泛的学术数据来源之一，其因子数据的构建方法参照Fama和French \cite{fama2015five} 的标准程序，在分组方式和加权方法上与国际文献保持一致。选择CSMAR因子数据而非自行从个股层面构建因子的主要考虑在于：第一，本文的研究重点在于VMP策略的有效性检验而非因子构建方法本身，使用公开标准化数据源有利于结果的可比性和可复现性；第二，CSMAR的因子数据已在Liu等人 \cite{liu2019size}、Hu等人 \cite{hu2019chinese}、李志冰等 \cite{li2017factor} 等多项中国市场因子研究中被广泛使用，其数据质量已经过学术界的反复验证。

\subsection{样本划分}

本文选择的全样本期间为2005年1月4日至2025年3月31日，共计4911个交易日。起始年份取2005年是因为此前A股尚处于股权分置改革前的特殊制度环境，非流通股与流通股并存，市值计算和因子排序的经济含义与改革后有本质差异，纳入会引入结构性断点。

全样本被划分为样本内（2010年1月至2019年12月）和样本外（2020年1月至2025年3月）两个子区间。前者覆盖了2015年急速牛熊转换和2018年持续下跌，后者包含新冠疫情冲击、注册制推行和北向资金大幅波动等新环境，用以检验策略的泛化能力。这一做法遵循Cederburg等人 \cite{cederburg2020performance} 建议的样本外检验范式，避免前视偏差。

另外选取三个危机子样本做压力测试：2015年股灾（6月至次年1月，突发流动性危机），2018年贸易战（全年，渐进式系统性下跌），以及2022年俄乌冲突与疫情封控（全年，外生冲击叠加政策不确定性）。三段时期的危机类型各不相同，用于考察VMP在不同情境下的保护效果。

\section{策略构建}

\subsection{波动率估计}

VMP策略的核心输入是对因子收益率条件方差的实时估计。本文采用滚动窗口已实现方差作为波动率测度。具体而言，对于每个因子$f$，在交易日$t$，基于过去$W$个交易日的日度收益率计算已实现方差：
\begin{equation}
  \hat{\sigma}_{f,t}^{2} = \sum_{i=1}^{W} r_{f,t-i}^{2}
  \label{eq:rv}
\end{equation}

其中$r_{f,t-i}$为因子$f$在$t-i$日的收益率，$W$为滚动窗口长度。此处使用收益率的平方和而非去均值后的平方和，与Moreira和Muir \cite{moreira2017volatility} 的设定一致——日度均值的数量级远小于标准差，去均值对估计精度几乎没有帮助。

基准窗口取$W = 21$个交易日（约一个自然月），这也是Moreira和Muir \cite{moreira2017volatility} 和Barroso和Santa-Clara \cite{barroso2015momentum} 的基准选择。稳健性分析中替换为$W = 63$（约一季度）和$W = 126$（约半年）。

\subsection{权重设定}

在获得条件方差估计后，VMP策略按照方差倒数的原则确定每日的投资权重。因子$f$在交易日$t$的VMP收益率定义为：
\begin{equation}
  r_{f,t}^{\text{VMP}} = \frac{c}{\hat{\sigma}_{f,t-1}^{2}} \cdot r_{f,t}
  \label{eq:vmp}
\end{equation}

其中$c$为缩放常数，$\hat{\sigma}_{f,t-1}^{2}$为$t-1$日计算的已实现方差，$r_{f,t}$为因子$f$在$t$日的实际收益率。直觉上说，前一日波动率越高，今日仓位越轻；波动率越低，仓位越重。策略使用$t-1$日信息决定$t$日权重，不存在前视偏差。

缩放常数$c$的设定目的是使VMP收益率序列的无条件标准差与原始因子收益率的无条件标准差保持一致：
\begin{equation}
  c = \frac{\text{Std}(r_{f})}{\text{Std}\left(\frac{r_{f}}{\hat{\sigma}_{f}^{2}}\right)}
  \label{eq:scaling}
\end{equation}

其中$\text{Std}(\cdot)$表示样本标准差。经过这一标准化后，VMP因子与原始因子在波动率水平上具有可比性，Spanning回归中的alpha就能直接反映同等风险下的增量收益，而不只是杠杆放大的结果。

\subsection{杠杆约束}

在实际执行中，公式(\ref{eq:vmp})可能产生过高的仓位权重。当因子进入极低波动率区间时，方差倒数$1/\hat{\sigma}_{f,t-1}^{2}$趋于极大值，对应的理论最优权重可能远超实际可融资水平。为确保策略的实践可行性并避免极端杠杆对结果的过度扭曲，本文对每日的投资权重施加上限约束：
\begin{equation}
  w_{f,t} = \min\left(\frac{c}{\hat{\sigma}_{f,t-1}^{2}},\ L_{\max}\right)
  \label{eq:leverage}
\end{equation}

其中$w_{f,t}$为因子$f$在$t$日的实际投资权重，$L_{\max} = 3$为最大杠杆倍数。选择3倍上限是因为中国市场融资融券的实际杠杆通常在1.5至2.5倍，3倍已略高于多数机构的实际融资能力，属于偏保守但不至于过度限制策略空间的设定。

\subsection{交易成本扣除}

VMP策略需要根据波动率变化逐日调整仓位权重，由此产生的交易成本是评估策略净收益时不可忽视的因素。本文按照每日权重变化的绝对值乘以单边交易费率计算交易成本：
\begin{equation}
  TC_{f,t} = |w_{f,t} - w_{f,t-1}| \times \kappa
  \label{eq:tc}
\end{equation}

其中$TC_{f,t}$为因子$f$在$t$日的交易成本，$w_{f,t}$和$w_{f,t-1}$分别为$t$日和$t-1$日的投资权重，$\kappa = 0.001$为单边交易费率（10个基点）。该费率综合考虑了佣金（机构端约万二至万三）、印花税分摊和市场冲击，属于较为保守的估计。

VMP因子的净收益率定义为：
\begin{equation}
  r_{f,t}^{\text{VMP,Net}} = w_{f,t} \cdot r_{f,t} - TC_{f,t}
  \label{eq:vmp_net}
\end{equation}

后文所有涉及VMP绩效的报告均基于扣除交易成本后的净收益率，以回应DeMiguel等人 \cite{demiguel2024transaction} 关于交易成本侵蚀VMP收益的批评。

\section{检验方法}

\subsection{绩效指标}

为全面评估VMP策略对各因子风险收益特征的影响，本文计算以下绩效指标。年化收益率定义为日均收益率乘以年交易日数$\bar{r} \times 252$。年化波动率定义为日收益率标准差乘以时间平方根$\text{Std}(r) \times \sqrt{252}$。夏普比率定义为年化收益率与年化波动率之比，衡量单位风险所获得的超额补偿。最大回撤定义为策略净值从历史高点到后续低点的最大跌幅：
\begin{equation}
  \text{MDD} = \max_t \left[\frac{\max_{s \leq t} V_s - V_t}{\max_{s \leq t} V_s}\right]
  \label{eq:mdd}
\end{equation}

其中$V_t$为策略在$t$日的累计净值。最大回撤刻画了投资者在最不利的入场时点可能承受的最大损失幅度，是实务中评估策略尾部风险容忍度的重要参考指标。此外，本文还报告了年化换手率，定义为每日权重变化绝对值之和的年化值$\sum |w_t - w_{t-1}| \times 252 / T$，用于衡量VMP策略的交易频率及其对交易成本假设的敏感程度。

\subsection{Spanning回归与Alpha检验}

VMP策略是否创造了超越原始因子的增量投资价值，需要通过正式的统计检验加以判断。本文采用Spanning回归框架，对每个因子$f$分别估计如下回归方程：
\begin{equation}
  r_{f,t}^{\text{VMP,Net}} = \alpha + \beta \cdot r_{f,t} + \varepsilon_t
  \label{eq:Spanning}
\end{equation}

该回归的经济含义在于：将VMP因子的净收益率对原始因子收益率进行投射，截距项$\alpha$衡量的是VMP组合中无法被原始因子线性复制的剩余收益部分。若$\alpha$在统计上显著为正，则说明一个持有原始因子的投资者通过转向VMP策略可以获得帕累托改善——在相同的因子暴露水平上获得更高的期望收益；若$\alpha$不显著或为负，则说明VMP策略未能提供增量价值，其收益完全可以被原始因子的适当线性组合所复制。$\beta$系数反映了VMP组合对原始因子的暴露程度，$R^2$则衡量原始因子对VMP收益变异的解释力。

由于日度收益率数据通常存在条件异方差和短期序列相关，OLS标准误可能偏低。为此，所有回归的$t$统计量均采用Newey-West \cite{newey1987simple} HAC标准误（滞后阶数设为5个交易日，参照Moreira和Muir \cite{moreira2017volatility} 的处理）。此外，本文对全样本、样本内、样本外以及三个危机子样本分别进行Spanning回归，以考察alpha在不同市场状态下的稳定性。


% ====================================================================
% 第五章：实证结果与分析
% ====================================================================
\chapter{实证结果与分析}

\section{描述性统计}

表~\ref{tab:desc_stats_raw}报告了六个原始因子日度收益率的描述性统计。全样本包含4911个交易日观测。从均值来看，市场因子（MKT\_RF）和规模因子（SMB）的日均收益率分别为$0.054\%$和$0.038\%$，对应的$t$统计量分别为$2.23$（$p=0.026$）和$3.29$（$p=0.001$），在统计上显著异于零，表明这两个因子在中国A股市场具有可靠的风险溢价。相比之下，价值因子（HML）、盈利因子（RMW）、投资因子（CMA）和动量因子（UMD）的日均收益率均不显著异于零，其中UMD的均值甚至为负（$-0.008\%$），印证了已有文献关于动量因子在中国市场无条件表现不佳的判断。

从分布特征来看，所有因子的Jarque-Bera统计量均在$1\%$水平下拒绝正态假设，这并不意外。比较值得注意的是SMB的超额峰度高达$8.44$，说明尖峰厚尾现象在规模因子上尤其突出。市场因子和规模因子的偏度为负——大跌出现的频率高于正态分布的预期——而盈利因子则呈正偏。这些厚尾特征正是VMP策略想要应对的风险来源：假如极端收益主要集中在高波动期，那么在高波动时降仓就有可能系统性地避开尾部损失。波动率方面，MKT\_RF日标准差最大（$1.72\%$），UMD次之（$1.05\%$），HML和CMA较低（$0.67\%$--$0.68\%$）。

\begin{table}[htbp]
\centering
\caption{原始因子日度收益率描述性统计}
\label{tab:desc_stats_raw}
\small
\begin{tabular}{l*{6}{r}}
\toprule
统计量 & MKT\_RF & SMB & HML & RMW & CMA & UMD \\
\midrule
观测数 & 4911 & 4911 & 4911 & 4911 & 4911 & 4911 \\
均值(\%) & 0.0544 & 0.0382 & 0.0077 & $-$0.0062 & 0.0002 & $-$0.0075 \\
标准差(\%) & 1.7245 & 0.8209 & 0.6816 & 0.8743 & 0.6695 & 1.0463 \\
偏度 & $-$0.5948 & $-$1.1147 & 0.2146 & 0.7937 & $-$0.0889 & $-$0.3503 \\
峰度 & 4.9203 & 8.4424 & 2.7141 & 6.3545 & 3.2901 & 2.1621 \\
最小值(\%) & $-$9.4780 & $-$8.0566 & $-$4.2669 & $-$5.9844 & $-$4.3705 & $-$6.0757 \\
最大值(\%) & 9.9204 & 6.1726 & 3.6150 & 8.5518 & 4.1618 & 5.0135 \\
\midrule
JB统计量 & 5328.99 & 15859.19 & 1569.53 & 8922.47 & 2256.99 & 1073.71 \\
$t$统计量 & 2.2331 & 3.2910 & 0.8017 & $-$0.5029 & 0.0192 & $-$0.5061 \\
$p$值 & 0.0256 & 0.0010 & 0.4228 & 0.6151 & 0.9847 & 0.6128 \\
\bottomrule
\end{tabular}

\vspace{0.5em}
\begin{minipage}{0.95\textwidth}
\footnotesize
\textit{注}：本表报告了2005年1月至2025年3月期间六个因子日度收益率的描述性统计。均值、标准差、最小值和最大值均以百分比表示。偏度和峰度为样本偏度和超额峰度。JB统计量为Jarque-Bera正态性检验统计量。$t$统计量检验均值是否显著异于零。
\end{minipage}
\end{table}

表~\ref{tab:desc_stats_vmp}报告了VMP策略管理后各因子日度净收益率的描述性统计（已扣除交易成本）。与原始因子相比，VMP因子在分布特征上呈现出若干值得关注的变化。第一，VMP显著增大了所有因子的超额峰度，尤其是MKT$^{\text{VMP}}$的峰度从$4.92$升至$18.14$。这一现象的成因在于VMP的动态杠杆机制：在低波动时期策略加大仓位暴露，一旦此时出现意外的极端收益，其幅度会被杠杆效应放大，从而在收益分布的尾部制造出更极端的观测值。第二，MKT和SMB经VMP管理后均值大幅下降——MKT\_RF从$0.054\%$降至$0.014\%$，SMB从$0.038\%$降至$0.005\%$，且$t$统计量均不再显著。这表明对于这两个本身具有显著正溢价的因子，VMP的动态调仓反而损害了其收益获取能力。第三，UMD$^{\text{VMP}}$的日均收益率从原始的$-0.008\%$转为正值$0.027\%$，对应$t$统计量为$1.96$（$p=0.050$），处于$5\%$显著性的临界值上，初步暗示VMP策略可能改善了动量因子的收益分布。

\begin{table}[htbp]
\centering
\caption{VMP因子日度净收益率描述性统计}
\label{tab:desc_stats_vmp}
\small
\begin{tabular}{l*{6}{r}}
\toprule
统计量 & MKT$^{\text{VMP}}$ & SMB$^{\text{VMP}}$ & HML$^{\text{VMP}}$ & RMW$^{\text{VMP}}$ & CMA$^{\text{VMP}}$ & UMD$^{\text{VMP}}$ \\
\midrule
观测数 & 4911 & 4911 & 4911 & 4911 & 4911 & 4911 \\
均值(\%) & 0.0140 & 0.0052 & 0.0071 & 0.0044 & $-$0.0082 & 0.0268 \\
标准差(\%) & 1.6088 & 0.7608 & 0.6448 & 0.8377 & 0.6202 & 0.9681 \\
偏度 & $-$1.5118 & $-$1.5135 & 0.5542 & 0.9963 & $-$0.2132 & 0.2026 \\
峰度 & 18.1367 & 9.2536 & 6.5074 & 6.3161 & 5.2077 & 4.7495 \\
最小值(\%) & $-$25.7072 & $-$6.9806 & $-$3.2128 & $-$4.2390 & $-$5.4778 & $-$5.6502 \\
最大值(\%) & 10.6950 & 3.0786 & 7.4657 & 6.7512 & 4.7697 & 6.6218 \\
\midrule
JB统计量 & 70328.74 & 19718.61 & 9062.50 & 9123.34 & 5677.30 & 4724.73 \\
$t$统计量 & 0.6134 & 0.4827 & 0.7797 & 0.3729 & $-$0.9374 & 1.9616 \\
$p$值 & 0.5397 & 0.6293 & 0.4356 & 0.7092 & 0.3486 & 0.0499 \\
\bottomrule
\end{tabular}

\vspace{0.5em}
\begin{minipage}{0.95\textwidth}
\footnotesize
\textit{注}：本表报告了波动率管理后各因子日度净收益率的描述性统计。VMP策略施加3倍最大杠杆约束，交易成本按每日权重变化幅度扣除单边万分之十计算。
\end{minipage}
\end{table}

表~\ref{tab:corr_matrix}报告了六个原始因子之间的Pearson相关系数矩阵。其中较为突出的特征是SMB与RMW之间高达$-0.83$的强负相关。此结果与Hu等人\cite{hu2019chinese}的研究发现相一致：在中国市场之中，小盘股盈利能力往往较弱，规模溢价与盈利溢价之间存在天然之对冲关系。此高相关性，也意味着于多因子框架之下同时纳入SMB与RMW之时，需要注意多重共线性对于回归系数估计精度的影响。此外，CMA与SMB（$0.50$）、CMA与HML（$0.49$）之间存在较高的正相关关系，而CMA与RMW之间则为$-0.50$的负相关，此结果反映了价值、投资与盈利因子于中国市场中的复杂联动关系。与上述因子间的高相关性形成对比的是，UMD与其他五个因子相关系数的绝对值，均不超过$0.21$，这表明动量因子提供了相对独立的信息来源。此种低相关性特征，为后文之Spanning回归提供了较为重要的背景：VMP对动量因子的改善效果是不太可能被其他因子所解释的，这是因为动量因子本身就处于因子空间的相对独立维度上。


\begin{table}[htbp]
\centering
\caption{因子收益率Pearson相关系数矩阵}
\label{tab:corr_matrix}
\small
\begin{tabular}{l*{6}{r}}
\toprule
 & MKT\_RF & SMB & HML & RMW & CMA & UMD \\
\midrule
MKT\_RF & 1.0000 &  &  &  &  &  \\
SMB & 0.1912 & 1.0000 &  &  &  &  \\
HML & $-$0.0389 & $-$0.1053 & 1.0000 &  &  &  \\
RMW & $-$0.2800 & $-$0.8277 & 0.0370 & 1.0000 &  &  \\
CMA & $-$0.0663 & 0.4983 & 0.4866 & $-$0.5048 & 1.0000 &  \\
UMD & 0.0464 & $-$0.1652 & $-$0.1732 & 0.1400 & $-$0.2064 & 1.0000 \\
\bottomrule
\end{tabular}

\vspace{0.5em}
\begin{minipage}{0.95\textwidth}
\footnotesize
\textit{注}：本表报告了六个因子日度收益率之间的Pearson线性相关系数。
\end{minipage}
\end{table}

\section{各因子VMP表现概览}

表~\ref{tab:perf_summary}报告了全样本（2005--2025年）下六个因子的原始收益与VMP净收益的核心绩效指标对比。一眼看过去就能发现，VMP对不同因子的影响方向差异很大，完全不是“普遍有效”的图景。

动量因子是唯一在收益、夏普比率和最大回撤三个维度上都明确受益的因子：年化收益从$-1.9\%$变成$+6.8\%$，夏普比率从$-0.11$到$0.44$，最大回撤从$77.1\%$收窄至$55.9\%$。而市场因子和规模因子则明显受损——MKT的夏普比率从$0.50$跌到$0.14$，SMB从$0.74$跌到$0.11$，两者的最大回撤在VMP介入后不降反升。价值、盈利和投资三个因子前后变化不大，基本在原地踏步。

\begin{table}[htbp]
  \centering
  \caption{各因子原始收益与VMP净收益绩效对比（全样本，21天窗口）}
  \label{tab:perf_summary}
  \begin{tabular}{lccccccc}
    \toprule
    因子 & 类型 & 年化收益 & 年化波动 & 夏普比率 & 最大回撤 \\
    \midrule
    \multirow{2}{*}{MKT} & 原始 & 13.7\% & 27.4\% & 0.50 & 73.8\% \\
                          & VMP净值 & 3.5\% & 25.5\% & 0.14 & 81.3\% \\
    \midrule
    \multirow{2}{*}{SMB} & 原始 & 9.6\% & 13.0\% & 0.74 & 49.3\% \\
                          & VMP净值 & 1.3\% & 12.1\% & 0.11 & 73.7\% \\
    \midrule
    \multirow{2}{*}{HML} & 原始 & 1.9\% & 10.8\% & 0.18 & 40.5\% \\
                          & VMP净值 & 1.8\% & 10.2\% & 0.17 & 47.7\% \\
    \midrule
    \multirow{2}{*}{RMW} & 原始 & $-1.6\%$ & 13.9\% & $-0.11$ & 67.8\% \\
                          & VMP净值 & 1.1\% & 13.3\% & 0.08 & 71.8\% \\
    \midrule
    \multirow{2}{*}{CMA} & 原始 & 0.05\% & 10.6\% & 0.00 & 57.7\% \\
                          & VMP净值 & $-2.1\%$ & 9.8\% & $-0.21$ & 54.9\% \\
    \midrule
    \multirow{2}{*}{UMD} & 原始 & $-1.9\%$ & 16.6\% & $-0.11$ & 77.1\% \\
                          & VMP净值 & 6.8\% & 15.4\% & 0.44 & 55.9\% \\
    \bottomrule
  \end{tabular}
\end{table}

\section{Spanning回归结果}

表~\ref{tab:Spanning_nw}报告了Spanning回归的结果，alpha是否显著为正是核心判断标准。

六个因子里，动量的alpha最突出：年化$+8.1\%$，$t = 3.27$（$p = 0.001$），在$1\%$水平下显著。意思是，如果一个投资者本来持有原始动量因子，换成VMP版本后，在保持同等因子暴露的情况下每年能多拿大约$8.1\%$的收益——这就是假设一所预测的效果。

规模因子的方向完全相反：alpha为$-5.4\%$（$t = -2.36$，$p = 0.018$），统计上确认了VMP对SMB的损害。市场因子alpha也是负的（$-5.2\%$），但$p = 0.261$，还不够显著，只能说方向与预期一致。HML、RMW、CMA三个因子的alpha都很小且远不显著，VMP对它们基本没有什么影响。
就$R^2$而言，各因子在$0.47$至$0.59$之间，原始因子大致能解释VMP因子一半到六成的收益变异。$\beta$都在$0.64$--$0.73$之间，均明显小于$1$，这反映了VMP在高波动期减仓导致的平均暴露下降。

\begin{table}[htbp]
  \centering
  \caption{Spanning回归结果（全样本，21天窗口，Newey-West修正）}
  \label{tab:Spanning_nw}
  \begin{tabular}{lcccccc}
    \toprule
    因子 & 年化Alpha & NW $t$值 & $p$值 & Beta & $R^2$ & 显著性 \\
    \midrule
    MKT   & $-5.2\%$ & $-1.13$ & 0.261 & 0.64 & 0.47 & 不显著 \\
    SMB   & $-5.4\%$ & $-2.36$ & 0.018 & 0.70 & 0.57 & 显著（负向）\\
    HML   & $+0.4\%$ & $+0.24$ & 0.810 & 0.71 & 0.56 & 不显著 \\
    RMW   & $+2.3\%$ & $+1.00$ & 0.317 & 0.73 & 0.58 & 不显著 \\
    CMA   & $-2.1\%$ & $-1.26$ & 0.208 & 0.70 & 0.58 & 不显著 \\
    UMD   & $+8.1\%$ & $+3.27$ & 0.001 & 0.71 & 0.59 & 显著（$1\%$水平） \\
    \bottomrule
  \end{tabular}
\end{table}

\section{2005-2025期间累积收益率曲线的对比分析}

\subsection{VMP对动量因子的改善效果}

图4-1为动量因子原始版本和VMP管理版本的累积收益率对比图。从图中可以清楚地看出，原始UMD的累积净值在2005年到2025年期间一直呈下降趋势，这与已有文献中关于中国A股市场动量效应失效的结论相吻合。尤其在2015年6月到7月的股灾期间，原始动量策略累积净值损失达到接近$50\%$，之后再也没有达到前期的高点。相反，VMP管理后的动量因子呈上升趋势。2015年股灾之前，原始动量组合因为市场突然反转而造成重大亏损的时候，VMP版本通过识别前期波动率急剧上升而提前大幅度降低了仓位暴露，从而有效地规避了这次极端损失事件。该机制同Barroso和Santa-Clara \cite{barroso2015momentum} 在美股动量策略里发现的保护作用十分相似，波动率管理可以准确地在动量崩盘发生之前把仓位降下来，因为动量崩盘总是事前波动率明显上升。


\begin{figure}[htbp]
  \centering
  \includegraphics[width=0.9\textwidth]{umd_Advanced_Cumulative_Return.png}
  \caption{动量因子（UMD）累积收益曲线对比：原始因子与VMP管理后净收益}
  \label{fig:umd_cumret}
\end{figure}

图~\ref{fig:umd_cumret}展示了动量因子的原始版本与VMP管理版本之间的累积收益率对比情况。从图中可以观察到，原始UMD的累积净值，在2005年至2025年间，整体呈现出下行趋势，此结果与已有文献关于中国A股市场动量效应失效的结论具有一定的一致性。尤其在2015年6月至7月的股灾期间，原始动量策略遭受了接近$50\%$的累积净值损失，且此后未能恢复至前期的高点。

相比之下，经VMP管理后的动量因子，呈现出相对稳定之上行趋势。这2015年股灾前后的对比尤为鲜明：当原始动量组合由于市场急剧反转而遭受重大亏损之时，VMP版本经由识别到前期波动率的急剧攀升，进而提前大幅降低了仓位暴露，由此有效回避了此一极端损失事件。此一机制，呼应了Barroso和Santa-Clara\cite{barroso2015momentum}的美股动量策略中所观察到的保护效果：由于动量崩盘通常伴随着事前波动率之显著升高，波动率管理能够于动量崩盘发生之前较为精确地降低仓位。

关于VMP在动量因子上有效性的经济解释，本文认为，其核心在于：中国A股市场的动量效应，并非完全不存在，而是于其收益分布的正中心部分包含了温和但持续之正动量信号。然此信号被少数高波动之极端反转事件所淹没。A股市场频繁出现的政策驱动型急速反转（例如2015年杠杆牛市后的流动性踩踏、2008年四万亿刺激后的风格切换），使得动量策略的输家组合在市场反转初期产生了远超赢家组合的涨幅，进而对动量策略造成了毁灭性的打击。由于此类极端反转事件时间上往往伴随着波动率的急剧升高，VMP经由在高波动时期自动降低仓位，能过系统性地规避了此类尾部损失事件。Daniel和Moskowitz\cite{daniel2016momentum}关于动量崩盘的研究也为这一解释提供了来自理论层面的支撑：其指出动量策略的尾部风险具有可预测性，而波动率恰恰是预测动量崩盘的有效信号之一。

\subsection{VMP对市场因子和规模因子的负面影响}

图~\ref{fig:mkt_cumret}与图~\ref{fig:smb_cumret}是市场因子、规模因子的累积收益率对比图。两个因子显示出比较接近的模式，也即原始因子的累积收益长期是上升的，VMP管理之后的版本在收益累积速度上明显慢于原始版本。

对于该现象的解释要回到VMP策略运作机制本身。市场因子和规模因子在A股市场上有着比较明显的正风险溢价（如表~\ref{tab:desc_stats_raw}所示，两者的$t$统计量分别为$2.23$和$3.29$），并且这些因子的高收益期常常和高波动时期之间存在正相关关系。就A股市场而言，其大幅上涨一般会出现在前期波动率急剧上升之后的市场反弹阶段，2008年底到2009年初、2014年底到2015年初、2020年7月的急速上涨，都是VMP策略因为识别到高波动而大幅降低仓位的时段。也就是说，VMP在市场反弹开始的时候就对因子暴露进行了系统性的削减，在这个时候一般会得到因子收益的大部分。Chi等人（2019）对于中国市场收益率和历史波动率之间存在正相关关系的发现给VMP在MKT、SMB两个因子上的表现不佳提供了一个实证方面的支持，当高波动预示着高预期收益的时候，根据波动率倒数调整仓位的策略就会造成收益损失。


\begin{figure}[htbp]
  \centering
  \includegraphics[width=0.85\textwidth]{mkt_rf_Advanced_Cumulative_Return.png}
  \caption{市场因子（MKT）累积收益曲线对比：原始因子与VMP管理后净收益}
  \label{fig:mkt_cumret}
\end{figure}

\begin{figure}[htbp]
  \centering
  \includegraphics[width=0.85\textwidth]{smb_Advanced_Cumulative_Return.png}
  \caption{规模因子（SMB）累积收益曲线对比：原始因子与VMP管理后净收益}
  \label{fig:smb_cumret}
\end{figure}

\subsection{VMP对其他因子的影响}
图~\ref{fig:hml_cumret}、图~\ref{fig:rmw_cumret}以及图~\ref{fig:cma_cumret}分别展示了价值因子、盈利因子与投资因子的累积收益率对比情况。

就价值因子（HML）而言，VMP管理前后的累积收益曲线走势高度相似，两者于整个样本期间保持近乎平行的轨迹，此结果与Spanning回归中alpha接近零（$+0.4\%$，$p = 0.810$）的结果，具有一定吻合性。价值因子在A股市场中的风险溢价本身较弱（日均收益$0.008\%$，$t = 0.80$），VMP波动率择时既未能从中提取出额外的收益信号，亦未对其造成显著损害。

盈利因子（RMW）的情况略有不同：原始因子于全样本期间的累积收益呈小幅下行趋势。VMP管理之后的收益略有改善，但是幅度极为有限（年化alpha仅$+2.3\%$，$p = 0.317$）。RMW于中国市场之特殊性，或在于其与SMB强负相关（$-0.83$）。此意味着盈利因子的收益波动，在很大程度上可以被规模因子所解释，因而VMP的独立调整空间受到了一定限制。

投资因子（CMA）在VMP管理后的表现，与原始因子基本持平，甚至略有恶化（年化alpha为$-2.1\%$，$p = 0.208$）。CMA在样本期间的总体表现接近零（日均收益$0.0002\%$），其属于A股市场中定价能力较弱的因子之一。对于本身不具有显著溢价的因子而言，VMP策略缺乏可以优化的收益空间，其交易成本反而可能对本已微薄的收益造成侵蚀。

\begin{figure}[H]
  \centering
  \includegraphics[width=0.85\textwidth]{hml_Advanced_Cumulative_Return.png}
  \caption{价值因子（HML）累积收益曲线对比：原始因子与VMP管理后净收益}
  \label{fig:hml_cumret}
\end{figure}

\begin{figure}[H]
  \centering
  \includegraphics[width=0.85\textwidth]{rmw_Advanced_Cumulative_Return.png}
  \caption{盈利因子（RMW）累积收益曲线对比：原始因子与VMP管理后净收益}
  \label{fig:rmw_cumret}
\end{figure}

\begin{figure}[H]
  \centering
  \includegraphics[width=0.85\textwidth]{cma_Advanced_Cumulative_Return.png}
  \caption{投资因子（CMA）累积收益曲线对比：原始因子与VMP管理后净收益}
  \label{fig:cma_cumret}
\end{figure}

\section{稳健性检验}

为验证上述核心发现——VMP策略显著改善动量因子的风险调整收益——是否依赖于特定的参数选择或样本区间，本节从三个维度进行稳健性检验。

\subsection{替换波动率估计窗口}

基准分析之中使用21天作为已实现方差之滚动窗口。表~\ref{tab:robust_window}将窗口替换为63天（约一季度）以及126天（约半年）进行稳健性检验。结果表明，随着窗口增大，VMP对动量因子的改善幅度有所递减。这或许是因为更长的窗口对短期波动率变化的敏感度有所降低，策略调仓频率亦相应减少。然而，UMD之alpha在三个窗口下，均在$1\%$水平显著（$p$值分别为$0.001$、$0.004$、$0.003$）。此表明结论不依赖于特定波动率估计参数。值得注意的是，63天窗口下$6.5\%$的年化alpha，与126天窗口下$6.1\%$的年化alpha相比非常接近，暗示着波动率信号中对VMP策略最为关键的低频成分在季度尺度上可能已基本被捕获。

\begin{table}[htbp]
  \centering
  \caption{不同波动率窗口下UMD因子VMP的Alpha检验（全样本）}
  \label{tab:robust_window}
  \begin{tabular}{lcccc}
    \toprule
    窗口 & 净夏普比率 & 年化Alpha & NW $t$值 & $p$值 \\
    \midrule
    21天（基准） & 0.44 & $+8.1\%$ & $+3.27$ & 0.001 \\
    63天 & 0.30 & $+6.5\%$ & $+2.91$ & 0.004 \\
    126天 & 0.24 & $+6.1\%$ & $+2.97$ & 0.003 \\
    \bottomrule
  \end{tabular}
\end{table}

\subsection{样本内与样本外检验}

表~\ref{tab:robust_insample}报告了样本内（2010年至2019年）与样本外（2020年至2025年）之UMD因子VMP表现情况。

样本内期间，VMP对动量因子的改善极为显著：alpha达到$10.9\%$（$p = 0.007$），VMP净夏普比率由$-0.04$提升至$0.61$。样本外期间，alpha的点估计值为$+4.3\%$，但是其统计显著性有所下降（$p = 0.242$）。此一结果需加以更审慎的解读：就经济意义而言，样本外的VMP净夏普比率（$0.48$），仍然为原始因子（$0.20$）的两倍以上，说明了VMP的风险管理效用于样本外期间并未消失；就统计意义而言，显著性下降，一方面或与样本外期间较短（约5年）所导致的检验功效不足有关，另一方面亦可能反映了2020年以来A股量化交易拥挤度的上升对动量策略生态所产生之影响。Cederburg等人\cite{cederburg2020performance}亦发现，VMP之样本外表现通常弱于样本内，此一模式于本文的中国市场数据中，得到了一定程度的再现。

\begin{table}[htbp]
  \centering
  \caption{UMD因子VMP净收益的样本内外对比（21天窗口）}
  \label{tab:robust_insample}
  \begin{tabular}{lccccc}
    \toprule
    时期 & 原始夏普 & VMP净夏普 & 年化Alpha & NW $t$值 & $p$值 \\
    \midrule
    样本内 (2010--2019) & $-0.04$ & 0.61 & $+10.9\%$ & $+2.71$ & 0.007 \\
    样本外 (2020--2025) & $+0.20$ & 0.48 & $+4.3\%$ & $+1.17$ & 0.242 \\
    \bottomrule
  \end{tabular}
\end{table}

\subsection{危机子样本分析}

表~\ref{tab:robust_crisis}报告了三次重大危机期间UMD因子的表现对比，其目的在于检验VMP于不同类型的极端市场环境中，是否具有差异化之保护效果。

结果显示，在2015年股灾期间，原始动量因子遭受了较为严重之亏损（年化$-85.0\%$），此为动量因子于A股样本期间内较为严重之崩盘事件之一。VMP将年化亏损缩窄至$-60.9\%$，最大回撤从$46.0\%$降至$34.5\%$，保护幅度约为三分之一。虽然VMP未能完全规避损失，然考虑到2015年股灾的极端突发性质——市场于极短时间内由牛市顶部转入流动性危机，波动率信号的前瞻期或较为有限——此保护幅度已属不易。

2018年贸易战期间，VMP之保护效果最为显著。在市场持续下跌的渐进式熊市环境中，VMP版动量因子之年化收益率达到$+30.2\%$ ，明显高于原始动量的$+6.4\%$，最大回撤仅为$10.6\%$ 。此一结果可为以下推断提供初步证据：对于波动率逐步攀升的渐进式下跌行情，VMP策略能够较为持续地识别风险信号，并及时降低仓位，其保护效果优于突发性危机。

2022年俄乌冲突与疫情封控期间，VMP与原始因子之表现差异或较小（年化收益$-22.6\%$ vs $-24.7\%$，最大回撤$24.0\%$ vs $25.0\%$）。此时期的波动率结构不同于前两次危机：市场波动并非完全由单一方向的趋势性下跌所主导，而是在政策预期反复变化之中，呈现出高频震荡特征，VMP的月度尺度波动率信号，对于此种日度级别的高频切换，反应也许不够灵敏。


\begin{table}[htbp]
  \centering
  \caption{重大危机期间UMD因子表现对比（21天窗口）}
  \label{tab:robust_crisis}
  \begin{tabular}{lcccc}
    \toprule
    危机时期 & 原始年化收益 & VMP净年化收益 & 原始最大回撤 & VMP净最大回撤 \\
    \midrule
    2015年股灾 & $-85.0\%$ & $-60.9\%$ & 46.0\% & 34.5\% \\
    2018年贸易战 & $+6.4\%$ & $+30.2\%$ & 17.9\% & 10.6\% \\
    2022年俄乌/疫情 & $-24.7\%$ & $-22.6\%$ & 25.0\% & 24.0\% \\
    \bottomrule
  \end{tabular}
\end{table}

综合三个维度的稳健性检验，VMP对动量因子的改善效果在不同的波动率窗口、不同的样本期间以及不同类型的市场危机环境中均得到支持，表明该结论具有较好的稳健性。


% ====================================================================
% 第六章：结论与启示
% ====================================================================
\chapter{结论与启示}

\section{核心结论}

本文将Moreira和Muir \cite{moreira2017volatility} 提出的波动率管理投资组合（VMP）策略应用于中国A股市场的六个核心定价因子，施加了3倍杠杆约束并按单边10个基点扣除交易成本。回顾全文的分析结果，最值得关注的发现是VMP策略在不同因子上产生了方向截然相反的效果，而非如Moreira和Muir在美股中报告的那样普遍有益。

在动量因子上，VMP表现出了明确且稳健的正面作用。经过波动率管理之后，UMD的年化收益由负转正（$-1.9\% \rightarrow +6.8\%$），夏普比率达到$0.44$，Spanning回归给出的年化Alpha为$8.1\%$（$t = 3.27$）。无论将滚动窗口替换为63天还是126天，该Alpha均在$1\%$水平下显著。2015年股灾和2018年贸易战子样本中VMP也展现了不同程度的保护作用。笔者认为这一结果的合理解释是：A股动量策略的收益分布中，正向信号被少数极端反转事件严重干扰了，而这些反转事件恰好集中于波动率急剧攀升的时段——VMP在事前识别到高波动并自动降仓，相当于把这些毁灭性的尾部亏损过滤掉了。换言之，A股动量并非没有信号，只是信号被噪音淹没了。

不过，在市场因子和规模因子身上，情况完全反了过来。SMB的VMP版本Alpha为$-5.4\%$（$p = 0.018$），MKT的Alpha为$-5.2\%$（经济上亏损幅度相当，但$p = 0.261$，未达到常规显著性水平）。两个因子的夏普比率在VMP介入后均出现了大幅缩水。这背后的逻辑并不复杂：A股的市场溢价和规模溢价在高波动期反而更丰厚——恐慌抛售后的反弹、政策转向带动的小盘股行情，恰恰发生在VMP判断“风险过高”而减仓的时候。结果就是策略系统性地把赚钱最多的日子低配了。

至于剩下三个因子——价值（HML）、盈利（RMW）和投资（CMA），VMP的净效果基本为零。Alpha的点估计值在$-2.1\%$到$+2.3\%$之间浮动，$p$值无一低于$0.20$。这倒也在预期之内，毕竟这三个因子的溢价来源是企业基本面的横截面差异，跟市场波动率高不高关系不大，VMP的仓位调整对它们而言既帮不上忙，也造不成太大伤害。

\section{理论价值与学术贡献}

本文实证结果中最有意思的一点在于：同样是“按波动率倒数调仓”这一个操作，作用于不同因子时居然产生了方向完全相反的后果。这种非对称性并非偶然，它实际上揭示了不同因子溢价来源之间的本质差别。市场溢价和规模溢价带有较强的风险补偿属性，高波动时期恰恰是风险补偿集中到账的窗口——此时减仓等于主动放弃回报最丰厚的时段。动量溢价则不同，其高波动期对应的往往是毁灭性崩盘而非丰厚回报，减仓反而起到了“避险”的效果。价值、盈利、投资这三个因子的情况介于两者之间，波动率的高低跟它们的收益关系不紧密，所以VMP的介入效果也就不痛不痒。说到底，VMP并非对一切因子都有效的万能工具，它的效果方向取决于一个简单的前提：高波动时期该因子的预期收益是正还是负。

这里还有一个方法论层面的启示。对动量因子做无条件均值检验（直接看全样本均值和$t$值），得到的结论是“A股不存在动量效应”。但经过VMP过滤之后，潜藏的正向信号浮现出来了。这多少说明：当一个因子的收益分布中混杂了与信号方向相反的极端尾部事件时，无条件检验会把有效信号和极端噪音搅在一起，从而低估甚至否定因子的定价能力。波动率管理在这里扮演了一种简单的条件过滤器角色——它不是在“创造”新的收益来源，而是在去除掩盖信号的极端噪音。

就具体的学术贡献而言，本文的工作可以从几个角度来定位。Cederburg等人\cite{cederburg2020performance}用美股数据指出，一旦纳入现实交易约束，VMP的Alpha只在动量因子上稳健；DeMiguel等人\cite{demiguel2024multifactor}在多因子框架下进一步确认了这个判断；Schwarz\cite{schwarz2025international}把分析推广到国际市场，发现VMP有效性存在跨市场的异质性。本文的结果——六个因子里只有动量一个的Alpha是显著的——为上述判断添加了一个来自中国新兴市场的独立样本点，与现有的国际证据形成了互补。

另一方面，关于A股动量效应为何“不灵”的问题，已有文献（如Liu等人\cite{liu2019size}）多从市场微观结构和投资者行为角度来解释。本文提供了一个不同的切入点：经VMP过滤后UMD产生了显著的正Alpha，这至少说明，A股动量效应的缺失不能简单归结为“信号不存在”，更准确的表述或许是“信号存在但被极端尾部事件所掩盖”。这一解释与Daniel和Moskowitz\cite{daniel2016momentum}关于动量崩盘可预测性的理论相吻合，也与Barroso和Santa-Clara\cite{barroso2015momentum}的实证发现相呼应。当然，本文的证据尚不能完全排除其他替代性解释，但至少给“A股无动量”这一结论打上了一个条件限定。

\section{实践启示}

对于A股量化投资实践来说，本文结果的一个直接含义是：波动率管理不是一个可以不加区分就套用到所有因子上的“通用增强器”。如果投资者持有的是市场因子或规模因子敞口——这两个因子本身已经具备相当稳定的正溢价——那么维持固定仓位可能反而比根据波动率来回调整更好。盲目套用VMP的后果，从本文数据来看，是夏普比率被腰斩甚至更多。

动量策略的情况则不同。裸动量在A股的历史表现不好，这是学界的共识，但本文的结果表明，一旦嵌入波动率管理的框架，动量策略就变得有投资价值了——夏普比率从负值拉到$0.44$，这在单因子策略中已经算不错的水平。所以对于有意在A股使用动量信号的投资者，一个实操建议是：不要照搬美股那种固定仓位的动量策略模式，而应该把实时波动率当作仓位管理的核心输入，在波动率飙升时果断减仓。

此外，在多因子组合配置中，因子自身近期的波动率水平或许可以作为差异化配置权重的参考依据。对于那些历史上回撤较大、表现不够稳定的因子，波动率信号也许有助于平滑极端市场环境下的组合净值曲线。当然，这里仍需提醒的是，交易成本会随调仓频率上升而累积，实际应用中需要在响应速度和成本控制之间做权衡。

\section{研究局限与未来方向}
本研究存在若干局限，有待后续研究加以改进。一个比较突出的问题是VMP对动量因子alpha的显著性在2020年以后的样本外期间有所下降——虽然经济意义上夏普比率仍然翻倍，但$p$值未达到常规显著水平。这或许跟近年A股量化交易拥挤度上升有关，动量策略的生态环境正在发生变化，这个问题值得持续追踪。另外，本文只做了单因子层面的VMP检验，如果把视野拓展到多因子组合，VMP和其他择时方法（比如宏观状态变量择时、条件均值预测等）之间是否存在协同效应，也是一个值得探索的方向。还有一点是波动率的度量方式——本文用的是最基础的滚动窗口已实现方差，未来可以尝试用高频数据构建的已实现波动率或者GARCH族模型的条件方差来替代，看更精确的波动率测度能否进一步提升VMP的效果。

%---------------------------------------------------------------------
%	参考文献
%---------------------------------------------------------------------

\printbibliography

%---------------------------------------------------------------------
%	致谢

\begin{acknowledgement} 

转瞬间，本科生涯已接近尾声。初进校园时对未知未来的期许、忐忑、憧憬与畅想仿佛就在昨天。在这一段旅程里，每一步的成长都离不开老师们的指导以及同学们的帮助。

首先，我要向我的指导老师陈莹教授表达感谢。四年来印象最深的课程之一便是陈莹老师的公司金融学。在毕业论文撰写的过程中，从论文选题的阶段开始，陈莹教授便以其在学术上的判断力为我指明了研究的方向。在之后进行数据分析以及撰写论文的过程中，老师也给予了我很多帮助。

其次，我要感谢南京大学所提供的学术环境，我想，在南大遇到的老师们，很可能是我此生能接触到的最厉害的学者。此外，学校为我们提供学术资源平台，如CSMAR等金融数据资源的开放共享，使得本科生也能够接触到实证研究所需要的工具。

最后，我要感谢我的家人和朋友们，感谢父母对我的指导、关怀和包容，感谢挚友们的陪伴，每个或欢乐或焦虑或通宵奋战的夜晚，都此生难忘。

总觉来日方长，原来不知不觉已经要走完本科阶段这许许多多的来日，愿我们都能一直顺利、健康、幸福。
\end{acknowledgement}

%---------------------------------------------------------------------
%	附录
%---------------------------------------------------------------------

\appendix

\chapter{核心数据处理代码}

本文的数据处理和VMP策略构建均使用Python实现。核心的计算流程包括：

\begin{enumerate}
  \item 从CSMAR日度因子数据中读取MKT\_RF, SMB, HML, RMW, CMA, UMD六个因子的日度收益率。
  \item 使用滚动窗口（21/63/126天）计算已实现方差，并取一阶滞后。
  \item 根据方差倒数计算原始权重，施加3倍杠杆上限，扣除基于权重变化的交易成本。
  \item 使用OLS回归配合Newey-West（HAC）标准误检验VMP净收益的Alpha显著性。
\end{enumerate}

\section{VMP策略构建核心代码}

以下代码实现了带有杠杆约束和交易成本的VMP策略构建。输入为因子日度收益率的DataFrame，输出为VMP净收益率序列。

\begin{lstlisting}[style=pythonstyle, caption={VMP策略构建函数}]
import pandas as pd
import numpy as np

def run_vmp(data, factors, window=21, leverage_limit=3.0, tc_bps=10):
    """
    构建带有杠杆约束和交易成本的VMP策略
    参数:
        data: DataFrame, 含因子日度收益率
        factors: list, 因子名称列表
        window: int, 已实现方差的滚动窗口(交易日)
        leverage_limit: float, 最大杠杆倍数
        tc_bps: float, 单边交易成本(基点)
    """
    df = data[factors].copy()

    # 步骤1: 计算滚动已实现方差及其一阶滞后
    for f in factors:
        df[f'{f}_RV'] = df[f].rolling(window).apply(
            lambda x: (x**2).sum(), raw=True)
        df[f'{f}_RV_lag'] = df[f'{f}_RV'].shift(1)
    df.dropna(inplace=True)

    # 步骤2: 对每个因子构建VMP
    for f in factors:
        raw_weight = 1.0 / df[f'{f}_RV_lag']
        # 缩放常数c: 使VMP波动率与原始因子一致
        c = df[f].std() / (df[f] * raw_weight).std()
        # 施加杠杆上限约束
        weight = (c * raw_weight).clip(upper=leverage_limit)
        df[f'{f}_Weight'] = weight
        # VMP毛收益
        df[f'{f}_VMP'] = weight * df[f]
        # 扣除交易成本: TC_t = |w_t - w_{t-1}| * kappa
        tc = weight.diff().abs().fillna(0) * (tc_bps / 1e4)
        df[f'{f}_VMP_Net'] = df[f'{f}_VMP'] - tc

    return df
\end{lstlisting}

\section{Spanning回归与Newey-West检验代码}

以下代码对VMP净收益进行Spanning回归，使用Newey-West异方差自相关一致标准误检验Alpha的显著性。

\begin{lstlisting}[style=pythonstyle, caption={Spanning回归与Alpha检验}]
import statsmodels.api as sm

def Spanning_regression_nw(data, factors, max_lags=5):
    """
    Spanning回归: r^{VMP,Net}_{f,t} = alpha + beta * r_{f,t} + eps_t
    使用Newey-West HAC标准误修正
    """
    results = []
    for f in factors:
        y = data[f'{f}_VMP_Net']  # VMP净收益
        X = sm.add_constant(data[f])  # 原始因子+截距
        # OLS + Newey-West HAC (maxlags=5)
        model = sm.OLS(y, X).fit(
            cov_type='HAC', cov_kwds={'maxlags': max_lags})
        results.append({
            'Factor': f,
            'Alpha_Ann': model.params['const'] * 252,
            't_NW': model.tvalues['const'],
            'p_value': model.pvalues['const'],
            'Beta': model.params[f],
            'R2': model.rsquared
        })
    return pd.DataFrame(results)
\end{lstlisting}

\section{描述性统计计算代码}

以下代码计算因子收益率的描述性统计量，包括Jarque-Bera正态性检验和均值t检验。

\begin{lstlisting}[style=pythonstyle, caption={描述性统计量计算}]
from scipy import stats

def descriptive_stats(data, factors):
    """
    计算因子日度收益率的描述性统计
    包含: 均值、标准差、偏度、峰度、JB检验、t检验
    """
    results = {}
    for f in factors:
        x = data[f].dropna()
        # 基本统计量
        results[f] = {
            'N': len(x),
            'Mean(%)': x.mean() * 100,
            'Std(%)': x.std() * 100,
            'Skewness': x.skew(),
            'Kurtosis': x.kurtosis(),  # 超额峰度
            'Min(%)': x.min() * 100,
            'Max(%)': x.max() * 100,
            'Median(%)': x.median() * 100,
        }
        # Jarque-Bera正态性检验
        jb, jb_p = stats.jarque_bera(x)
        results[f]['JB_stat'] = jb
        # t检验: H0: mu = 0
        t, p = stats.ttest_1samp(x, 0)
        results[f]['t_stat'] = t
        results[f]['p_value'] = p
    return pd.DataFrame(results)
\end{lstlisting}

\end{document}
